%--------------------------------------------------------------------------------
%
%
%--------------------------------------------------------------------------------
\pagebreak
\unnumberedSection{PxCE - Purge Instruction or Data Cache Entries}

\begin{iPageDescEntry}{Syntax}
    \texttt{PICE [ RegX] ( RegB )} \\
    \texttt{PDCE [ RegX] ( RegB )}
\end{iPageDescEntry}

\begin{iPageDescEntry}{Format}
    The \texttt{PICE} instruction uses the following instruction format: \\

    \begin{flushleft}
    \drawInstrFormat
        {0,1,3,5,6.5,8.5,9.5,11.5,16}
        {\OpCodeGrpSYS, \OpCodePCA, 0, 6, RegB, 0, RegX, 0}
    \end{flushleft}
    
    The \texttt{PDCA} instruction uses the following instruction format. \\

    \begin{flushleft}
    \drawInstrFormat
        {0,1,3,5,6.5,8.5,9.5,11.5,16}
        {\OpCodeGrpSYS, \OpCodePCA, 0, 7, RegB, 0, RegX, 0}
    \end{flushleft}
\end{iPageDescEntry}

\begin{iPageDescEntry}{Description}
    The \texttt{PxCA} instructions purge a cache entry or a set of cache entries selected by the effective address.
    in \texttt{RegB}. This can be used by privileged code for cache management. Unlike \texttt{FxCA}, a purge operation may also invalidate dirty lines without writing them back.
   
\end{iPageDescEntry}

\begin{iPageDescEntry}{Operation}
    \texttt{ purgeEntriesFromCache( addOfs32( RegB, RegX ));}
\end{iPageDescEntry}

\begin{iPageDescEntry}{Exceptions}
    None.
\end{iPageDescEntry}

\begin{iPageDescEntry}{Notes}
    Typically only allowed in supervisor mode.
\end{iPageDescEntry}
